\documentclass[a4paper]{article}

\usepackage[portuguese]{babel}
\usepackage[utf8]{inputenc}
\usepackage{indentfirst}
\usepackage{graphicx}
\usepackage{verbatim}

\begin{document}

\setlength{\textwidth}{16cm}
\setlength{\textheight}{22cm}

\title{\Huge\textbf{Jogos de tabuleiro}\linebreak\linebreak\linebreak
\Large\textbf{Relatório 1ª Fase}\linebreak\linebreak
\includegraphics[height=6cm, width=7cm]{feup.pdf}\linebreak \linebreak
\Large{Mestrado Integrado em Engenharia Informática e Computação} \linebreak \linebreak
\Large{Base de dados}\linebreak
}

\author{\textbf{Grupo 601:}\\ Hugo Ari Rodrigues Drumond --- 201102900 \\  Ricardo Jorge Matos Figueiredo --- 201100687\\\linebreak\linebreak \\
 \\ Faculdade de Engenharia da Universidade do Porto \\ Rua Roberto Frias, 4200--65 Porto, Portugal \linebreak\linebreak\linebreak
\linebreak\linebreak\vspace{1cm}}
\maketitle
\thispagestyle{empty}

\newpage

\section*{Resumo}
%Decidimos elaborar uma base de dados para uma Empresa organizadora de \textit{Jogos de tabuleiro}.    \cite{creator}


%Descrever muito sumariamente (1-2 parágrafos) o trabalho que está a ser reportado

%\section{Introdução}

%Descrever os objectivos e motivação do trabalho.
%Todas as figuras devem ser referidas no texto. %\ref{fig:codigoFigura}

\section{Contexto}
Esta base de dados destina-se a um Salão de jogos que organiza jogos de Tabuleiro.  
%Para quem se destina, quem e que vai usar, 

%\clearpage
\section{Conceitos Principais}
%Identificar as classes e falar um pouco sobre elas. As óbvias não é necessário.

\section{Diagrama de classes UML}
%Incluir todas as restrições, digitalizar o diagrama ou fazer num programa qualquer. Incluir notas.


\clearpage
\addcontentsline{toc}{section}{Bibliografia}
\renewcommand\refname{Bibliografia}
\bibliographystyle{plain}
\bibliography{myRefs}

\end{document}
